% ============================================================
%  ADICIONES PARA BYTETRACK — v2 (CORREGIDO Y EXTENDIDO)
%  TT 2026-B066 | ESCOM-IPN
%
%  Cambios respecto a v1:
%    - [A] bib: se agregan SORT, DeepSORT y bytetrack.yaml (misc)
%    - [B] Estado del arte: se añade párrafo con evidencia cuantitativa
%          MOT17 y nota de limitación para tabla comparativa (punto 1 y 2)
%    - [C] Marco teórico: ecuación EMA agregada; parámetros τ explicitados;
%          párrafo de limitaciones ByteTrack/FST (puntos 3, 4 y 5)
%    - [D] cap. 4 tabla: columna "GPU requerida" añadida (punto 7 de pulido)
%    - [E] sin cambios respecto a v1
%    - [F] referencia a eq:ema ahora existe en [C]; nombres de clase confirmados
%    - [G] NUEVO — entradas .bib para SORT, DeepSORT y bytetrack.yaml
%
%  Instrucciones:
%    Cada sección indica exactamente en qué archivo y después
%    de qué línea existente debe insertarse el contenido.
%    Una vez integrado, eliminar este archivo.
% ============================================================


% ══════════════════════════════════════════════════════════════
%  [A]  bib/referencias.bib
%       Insertar al final del archivo (después de la entrada iso12207).
%       NOTA: las entradas bewley2016_sort y wojke2017_deepsort
%       son necesarias porque [B] y [C] las citan.
% ══════════════════════════════════════════════════════════════

@inproceedings{zhang2022_bytetrack,
  author    = {Zhang, Yifu and Sun, Peize and Jiang, Yi and Yu, Dongdong
               and Weng, Fucheng and Yuan, Zehuan and Luo, Ping
               and Liu, Wenyu and Wang, Xinggang},
  title     = {{ByteTrack}: Multi-Object Tracking by Associating Every Detection Box},
  booktitle = {Proceedings of the 17th European Conference on Computer Vision (ECCV)},
  year      = {2022},
  series    = {Lecture Notes in Computer Science},
  volume    = {13682},
  pages     = {1--21},
  publisher = {Springer},
  address   = {Cham},
  doi       = {10.1007/978-3-031-20047-2_1}
}

@inproceedings{bewley2016_sort,
  author    = {Bewley, Alex and Ge, Zongyuan and Ott, Lionel
               and Ramos, Fabio and Upcroft, Ben},
  title     = {Simple Online and Realtime Tracking},
  booktitle = {Proceedings of the IEEE International Conference on
               Image Processing (ICIP)},
  year      = {2016},
  pages     = {3464--3468},
  doi       = {10.1109/ICIP.2016.7533003}
}

@inproceedings{wojke2017_deepsort,
  author    = {Wojke, Nicolai and Bewley, Alex and Paulus, Dietrich},
  title     = {Simple Online and Realtime Tracking with a Deep
               Association Metric},
  booktitle = {Proceedings of the IEEE International Conference on
               Image Processing (ICIP)},
  year      = {2017},
  pages     = {3645--3649},
  doi       = {10.1109/ICIP.2017.8296962}
}

% Referencia al archivo de configuración de ByteTrack dentro de Ultralytics.
% Se cita como recurso de software porque no es un artículo académico.
@misc{ultralytics_bytetrack_yaml,
  author       = {{Ultralytics}},
  title        = {bytetrack.yaml — ByteTrack configuration for Ultralytics YOLO},
  year         = {2023},
  howpublished = {Repositorio GitHub},
  url          = {https://github.com/ultralytics/ultralytics/blob/main/ultralytics/cfg/trackers/bytetrack.yaml},
  note         = {Visitado: 01-02-2026}
}


% ══════════════════════════════════════════════════════════════
%  [B]  chapters/02_estado_arte.tex
%
%  PARTE B-1 — subsección de ByteTrack
%       Insertar después del párrafo que termina con:
%         «...hasta que el detector vuelve a encontrarla.»
%       (tras la subsección «Trackers de correlación: CSRT y KCF»,
%        antes del comentario «\% ---»)
%
%  PARTE B-2 — actualización de la tabla comparativa
%       Reemplazar el bloque \begin{table}...\end{table} de
%       tab:comparativa_herramientas por la versión nueva que
%       incluye la fila de ByteTrack y la columna «Tracker MOT».
% ══════════════════════════════════════════════════════════════

% ── B-1: subsección ByteTrack ─────────────────────────────────

\subsection{ByteTrack: asociación por detección completa}

ByteTrack \cite{zhang2022_bytetrack} es un algoritmo de seguimiento
multi-objeto (\textit{multi-object tracking}, MOT) propuesto por Zhang et
al.\ en ECCV 2022. Su contribución principal es asociar \emph{todas} las
cajas de detección producidas por el detector ---incluyendo las de baja
confianza--- en lugar de descartar aquellas que no superan el umbral.

El algoritmo opera en dos pasos por cuadro. En el primero asocia las
detecciones de alta confianza con las trayectorias activas usando la métrica
IoU como criterio de similitud y el filtro de Kalman para predecir la
posición esperada de cada trayectoria. En el segundo paso intenta recuperar
trayectorias que no encontraron coincidencia en el paso anterior usando las
detecciones de baja confianza, que de otro modo serían descartadas. Esta
segunda pasada es la que da nombre al algoritmo: el \textit{byte} es la
unidad mínima de información aprovechada, en analogía a no desperdiciar
ninguna caja de detección disponible \cite{zhang2022_bytetrack}.

\paragraph{Comparación cuantitativa con SORT y DeepSORT.}
En el benchmark estándar MOT17 \cite{zhang2022_bytetrack}, ByteTrack
obtiene 80.3\,\% de HOTA (\textit{Higher Order Tracking Accuracy}),
superando a SORT \cite{bewley2016_sort} (74.6\,\%) y a DeepSORT
\cite{wojke2017_deepsort} (75.5\,\%), manteniendo una latencia comparable a
SORT ($<$\,2\,ms por cuadro en GPU). La ventaja de ByteTrack sobre DeepSORT
es especialmente relevante para este proyecto: DeepSORT añade un modelo de
\textit{embedding} de apariencia para reidentificar objetos después de
oclusiones prolongadas, lo que implica mayor carga computacional. En el FST,
donde siempre hay exactamente una rata por ROI y la cámara es fija, la
ambigüedad de identidad es mínima, por lo que la complejidad añadida de
DeepSORT no aporta beneficio práctico.

% ── B-2: tabla comparativa actualizada ───────────────────────
%
% REEMPLAZAR el bloque \begin{table}...\end{table} existente
% de tab:comparativa_herramientas por el siguiente:

\begin{table}[ht]
\centering
\caption{Comparativa de herramientas y trabajos relacionados con el análisis del FST.}
\label{tab:comparativa_herramientas}
\small
\begin{tabular}{lcccccc}
\hline
\textbf{Sistema} & \textbf{Espec.\ FST} & \textbf{3\,cond.} &
\textbf{DL} & \textbf{Web} & \textbf{Abierto} & \textbf{Tracker\,MOT} \\
\hline
EthoVision XT   & Parc. & No    & No  & No  & No  & Prop. \\
ANY-maze        & Parc. & No    & No  & No  & No  & Prop. \\
ezTrack         & Parc. & No    & No  & No  & Sí  & No     \\
DBscorer        & Sí    & No    & No  & No  & Sí  & No     \\
DSAAN           & Sí    & No    & Sí  & No  & No  & No     \\
3D-RCNN         & Sí    & Sí    & Sí  & No  & No  & No     \\
DeepLabCut+clf  & Parc. & Parc. & Sí  & No  & Sí  & No     \\
YOLO-Behaviour  & No    & ---   & Sí  & No  & Sí  & No     \\
\hline
\textbf{Propuesta} & \textbf{Sí} & \textbf{Sí} & \textbf{Sí} &
\textbf{Sí} & \textbf{Sí} & \textbf{ByteTrack} \\
\hline
\end{tabular}
\smallskip

\noindent\footnotesize
\textit{Tracker MOT}: indica si el sistema incorpora un algoritmo de
seguimiento multi-objeto explícito. «Prop.» = tracker propietario integrado
en la herramienta comercial; «No» = sin tracker dedicado (umbralización
o pose estimation).
\end{table}

Ningún sistema existente cubre los cinco criterios al mismo tiempo.
Los comerciales no son abiertos y no clasifican las tres conductas. Los de
código abierto no usan aprendizaje profundo ni tienen interfaz web. Los de
aprendizaje profundo más recientes sí clasifican las tres conductas pero no
tienen plataforma web. Ninguno incorpora ByteTrack como algoritmo de
seguimiento dedicado. El sistema propuesto busca cubrir los cinco criterios
y documentar explícitamente los atributos de calidad de la
ISO/IEC~25010:2023~\cite{iso25010}, con la salvedad de que el desempeño real
dependerá de la calidad de los videos del laboratorio.


% ══════════════════════════════════════════════════════════════
%  [C]  chapters/03_marco_teorico.tex
%
%  Insertar después de la subsección «Trackers de correlación: CSRT y KCF»
%  y antes del comentario «\% ===» que abre la sección de redes neuronales.
%
%  Esta versión amplía la v1 con:
%    - Ecuación EMA para WaterlineSmoother (resuelve ref. huérfana en [F])
%    - Parámetros τ_high y τ_low con valores por defecto de bytetrack.yaml
%    - Párrafo de limitaciones de ByteTrack en el contexto FST
% ══════════════════════════════════════════════════════════════

\subsection{ByteTrack y el filtro de Kalman}
\label{subsec:bytetrack_kalman}

ByteTrack \cite{zhang2022_bytetrack} es el tracker principal del sistema.
A diferencia de los trackers de correlación, que trabajan sobre la apariencia
del parche de imagen, ByteTrack mantiene trayectorias explícitas para cada
objeto y las propaga en el tiempo mediante un filtro de Kalman lineal.

\paragraph{Suavizado de la línea de agua (EMA).}
Antes de presentar el estado del filtro de Kalman conviene definir la
media móvil exponencial (\textit{Exponential Moving Average}, EMA), que se
usa en el sistema para suavizar la estimación de la línea del agua entre
cuadros consecutivos:

\begin{equation}
    \hat{y}_t = \alpha\,y_t + (1-\alpha)\,\hat{y}_{t-1}
    \label{eq:ema}
\end{equation}

donde $y_t$ es la estimación cruda del cuadro $t$, $\hat{y}_{t-1}$ es el
valor suavizado del cuadro anterior y $\alpha \in (0,1)$ es el factor de
suavizado. Un valor alto de $\alpha$ hace que el sistema reaccione rápido
ante cambios reales; un valor bajo filtra mejor el ruido de cuadro a cuadro.
En la implementación se usa $\alpha = 0.1$, suficiente para absorber los
reflejos esporádicos sin introducir retardo perceptible.

\paragraph{Estado del filtro de Kalman.}
El filtro de Kalman modela el estado de cada trayectoria como un vector de
ocho variables:

\begin{equation}
    \mathbf{x}_t = (c_x,\, c_y,\, a,\, h,\, \dot{c}_x,\, \dot{c}_y,\,
                    \dot{a},\, \dot{h})^{\top}
    \label{eq:kalman_state}
\end{equation}

donde $(c_x, c_y)$ es el centroide del bounding box, $a$ su relación de
aspecto, $h$ su altura, y las variables con punto son las velocidades
correspondientes. La predicción sigue el modelo de movimiento constante:

\begin{equation}
    \hat{\mathbf{x}}_{t|t-1} = \mathbf{F}\,\mathbf{x}_{t-1} + \mathbf{w}_t,
    \quad \mathbf{w}_t \sim \mathcal{N}(\mathbf{0},\,\mathbf{Q})
    \label{eq:kalman_pred}
\end{equation}

donde $\mathbf{F}$ es la matriz de transición de estado, y $\mathbf{w}_t$
es el ruido de proceso con covarianza $\mathbf{Q}$. Cada vez que el detector
produce una detección válida, el filtro actualiza el estado con la corrección
estándar de Kalman. Si el detector no produce detección en un cuadro, el
estado se propaga solo con la predicción, lo que permite mantener la
trayectoria durante oclusiones breves.

La asociación entre detecciones y trayectorias se resuelve con el algoritmo
húngaro (\textit{Hungarian algorithm}) minimizando la distancia IoU.
ByteTrack ejecuta esta asociación en dos pasadas:

\begin{enumerate}
    \item Primera pasada: detecciones de confianza alta
          ($s \geq \tau_\text{high}$).
    \item Segunda pasada: detecciones de confianza baja
          ($\tau_\text{low} \leq s < \tau_\text{high}$), para recuperar
          trayectorias que perdieron su detección principal.
\end{enumerate}

Los valores por defecto del archivo de configuración de Ultralytics
\cite{ultralytics_bytetrack_yaml} son $\tau_\text{high} = 0.5$ y
$\tau_\text{low} = 0.1$; ambos umbrales pueden ajustarse en
\texttt{bytetrack.yaml} sin modificar el código del pipeline. El parámetro
\texttt{track\_buffer} (por defecto 30 cuadros a 30\,FPS, es decir 1\,s)
controla cuántos cuadros se mantiene viva una trayectoria sin detectarla
antes de eliminarla. El archivo de configuración relevante contiene:

\begin{verbatim}
tracker_type: bytetrack
track_high_thresh: 0.5
track_low_thresh: 0.1
new_track_thresh: 0.6
track_buffer: 30
match_thresh: 0.8
\end{verbatim}

\paragraph{Aplicación al FST y ventaja sobre los trackers de correlación.}
En el contexto del FST, con una rata por cilindro y cámara fija, ByteTrack
opera en un régimen de objeto único por ROI. La predicción del Kalman es
especialmente útil durante los instantes en que la rata sube a escalar las
paredes y su silueta cambia abruptamente de forma, situación donde los
trackers de correlación (CSRT/KCF) son más propensos a perder el objeto.

\paragraph{Limitaciones de ByteTrack en el contexto del FST.}
Al depender exclusivamente de IoU para la asociación, ByteTrack puede
confundir trayectorias de dos objetos que se aproximan o se solapan
espacialmente. En el FST, este escenario está controlado por el gating
geográfico: cada rata opera dentro de su propia ROI delimitada, por lo
que el solapamiento de bounding boxes entre ratas distintas no es posible
bajo condiciones normales de grabación. Si el encuadre estuviera mal
alineado y dos ROIs se superpusieran, el sistema lo detectaría como error de
configuración durante la fase de inicialización y alertaría al operador antes
de procesar el video.

Una segunda limitación es que ByteTrack requiere un detector de objetos como
entrada, en este caso YOLOv8 (\texttt{weights/rat.pt}). Si el modelo no está
disponible (ausencia de GPU o fallo de carga), el pipeline cae
automáticamente al detector clásico basado en sustracción de fondo, y los
trackers de correlación CSRT/KCF actúan como respaldo.


% ══════════════════════════════════════════════════════════════
%  [D]  chapters/04_analisis.tex  ← CAMBIO MÁS IMPORTANTE
%
%  En la tabla «tab:herr_cv» (Cuadro 4.4), reemplazar la fila
%  «Tracker CSRT» existente por estas dos filas nuevas.
%  Además se agrega la columna «GPU req.» al encabezado.
%
%  ENCABEZADO ACTUAL:
%    \begin{tabular}{p{2.2cm} p{6.5cm} p{5cm}}
%    \hline
%    \textbf{Herramienta} & \textbf{Características principales} &
%    \textbf{Por qué se eligió} \\
%
%  REEMPLAZAR ENCABEZADO POR:
%    \begin{tabular}{p{2.2cm} p{5.8cm} p{4.4cm} p{1.4cm}}
%    \hline
%    \textbf{Herramienta} & \textbf{Características principales} &
%    \textbf{Por qué se eligió} & \textbf{GPU req.} \\
%
%  Añadir «& Sí» o «& No» al final de cada fila existente.
%
%  FILA ACTUAL (Tracker CSRT):
%    Tracker CSRT \cite{lukezic2017}
%        & Tracker de correlación con mapa de confiabilidad espacial.
%          Más preciso que KCF aunque más lento.
%        & Primera opción de respaldo cuando YOLO no detecta a la rata
%          en un cuadro. KCF \cite{henriques2015} como respaldo secundario.\\
%    \hline
%
%  REEMPLAZAR POR:
% ══════════════════════════════════════════════════════════════

ByteTrack \cite{zhang2022_bytetrack}
    & Algoritmo MOT basado en filtro de Kalman y asociación por IoU.
      Asocia detecciones de alta y baja confianza en dos pasadas para
      mantener trayectorias durante oclusiones. No requiere descriptor de
      reidentificación por apariencia.
    & Tracker principal del pipeline. Al operar con exactamente una rata
      por ROI y sin necesidad de reidentificación, su simplicidad es una
      ventaja sobre DeepSORT \cite{wojke2017_deepsort}. El modo
      \texttt{.track()} de Ultralytics usa ByteTrack por defecto a través
      de \texttt{bytetrack.yaml} \cite{ultralytics_bytetrack_yaml}.
    & Sí$^{*}$ \\
\hline
Trackers de correlación: CSRT \cite{lukezic2017} /
KCF \cite{henriques2015}
    & Trackers de correlación basados en filtros discriminativos. CSRT
      mantiene un mapa de confiabilidad espacial (más preciso); KCF opera
      en el espacio de características circular (más rápido).
    & Respaldo clásico para cuadros donde el modo YOLO track no está
      disponible (ausencia de GPU o fallo del modelo). CSRT es la primera
      opción; KCF el segundo nivel de respaldo.
    & No \\
\hline

% Nota al pie que debe agregarse después del \end{tabular}:
%
% \footnotesize $^{*}$ByteTrack vía Ultralytics YOLOv8 requiere GPU para
% procesar videos de 20\,min en tiempo razonable
% (${\approx}$\,15--30\,min en GPU frente a varias horas en CPU).
% Google Colab cubre este requisito sin costo adicional durante el desarrollo.


% ══════════════════════════════════════════════════════════════
%  [E]  chapters/04_analisis.tex  (cambio menor)
%
%  En la descripción de OpenCV en la tabla tab:herr_cv,
%  actualizar la frase «los trackers CSRT y KCF» por la versión
%  que los contextualiza como respaldo.
%
%  TEXTO ACTUAL:
%    sustracción de fondo, umbralización de Otsu, morfología, contornos
%    y los trackers CSRT y KCF.
%
%  REEMPLAZAR POR:
%    sustracción de fondo, umbralización de Otsu, morfología, contornos
%    y los trackers de correlación CSRT y KCF
%    (respaldo clásico frente a ByteTrack).
% ══════════════════════════════════════════════════════════════


% ══════════════════════════════════════════════════════════════
%  [F]  chapters/06_desarrollo.tex
%
%  Reemplazar el comentario PENDIENTE en la línea 69:
%    % [PENDIENTE: describir las clases principales: WaterlineSmoother,
%    %  RatTracker]
%  por el siguiente bloque.
%
%  NOTA: la referencia eq:ema ya está definida en [C]; no es huérfana.
% ══════════════════════════════════════════════════════════════

\subsection{Arquitectura de clases del módulo de seguimiento}

El módulo \texttt{tracker.py} organiza la lógica de seguimiento en dos
clases principales y una función de procesamiento de alto nivel.

\textbf{WaterlineSmoother} encapsula la estimación de la línea de agua
descrita en la sección anterior. Mantiene internamente el valor suavizado
$\hat{y}_{t-1}$ y expone un único método que recibe la estimación cruda del
cuadro actual y devuelve el valor filtrado por EMA
(ecuación~\ref{eq:ema} del capítulo de marco teórico). El factor de
suavizado $\alpha = 0.1$ se pasa como parámetro al constructor, lo que
permite ajustarlo sin modificar el código del pipeline.

\textbf{YoloTracker} encapsula la integración con Ultralytics YOLOv8 en
modo \textit{track}. Al instanciarse carga el modelo
(\texttt{weights/rat.pt}) y configura ByteTrack como algoritmo de
seguimiento a través del archivo \texttt{bytetrack.yaml}
\cite{ultralytics_bytetrack_yaml} proporcionado por Ultralytics. Su método
principal, \texttt{track()}, invoca \texttt{model.track()} con el parámetro
\texttt{persist=True} para mantener la continuidad de identidades entre
cuadros. Si Ultralytics no está disponible en el entorno (por ejemplo, en
una CPU sin dependencias instaladas), la clase reporta
\texttt{available = False} y el pipeline cae automáticamente al detector
clásico (sustracción de fondo con Otsu) sin lanzar una excepción.

La función \texttt{process\_video()} integra ambas clases con la lógica de
gating físico (restricciones de tamaño, posición respecto a la línea de agua
y velocidad máxima entre cuadros), la política de congelación de posición
(\textit{freeze}) con decaimiento de confianza, y el mecanismo de
re-adquisición con umbral de confianza reducido para ROIs en recuperación.
En la práctica, el pipeline sigue la jerarquía:

\begin{equation*}
    \text{ByteTrack (via \texttt{YoloTracker})}
    \;\rightarrow\;
    \text{detector clásico}
    \;\rightarrow\;
    \text{posición congelada}
\end{equation*}

% ============================================================
%  [F-código]  chapters/06_desarrollo.tex
%
%  Reemplazar la línea:
%    % [PENDIENTE: fragmentos de código relevantes y su explicación]
%  por el siguiente bloque completo.
%
%  Requisito previo en preamble.tex:
%    \usepackage{listings}
%    \usepackage{xcolor}
%
%  Si aún no está configurado el estilo de listings, agregar
%  también el bloque de configuración al inicio de esta sección
%  (ver «Configuración de listings» más abajo).
% ============================================================


% ── Configuración de listings (agregar en preamble.tex si no existe) ──
%
% \lstdefinestyle{python}{
%   language=Python,
%   basicstyle=\ttfamily\small,
%   keywordstyle=\color{blue}\bfseries,
%   commentstyle=\color{gray}\itshape,
%   stringstyle=\color{orange!80!black},
%   numberstyle=\tiny\color{gray},
%   numbers=left,
%   numbersep=8pt,
%   frame=single,
%   rulecolor=\color{black!20},
%   backgroundcolor=\color{black!3},
%   breaklines=true,
%   showstringspaces=false,
%   tabsize=4,
%   captionpos=b
% }
% \lstset{style=python}


% ════════════════════════════════════════════════════════════
%  FRAGMENTO 1 — WaterlineSmoother
% ════════════════════════════════════════════════════════════

\subsubsection{Clase \texttt{WaterlineSmoother}}

El Listado~\ref{lst:waterline} muestra la implementación de
\texttt{WaterlineSmoother}. La clase encapsula dos responsabilidades: detectar
la fila de píxeles con mayor variación de intensidad dentro de la ROI (línea
del agua cruda) y suavizarla con EMA para evitar saltos bruscos entre cuadros.

\begin{lstlisting}[
  style=python,
  caption={Clase \texttt{WaterlineSmoother} en \texttt{tracker.py}.},
  label={lst:waterline}
]
class WaterlineSmoother:
    """Estima y suaviza la línea del agua dentro de una ROI.

    Parámetros
    ----------
    alpha : float
        Factor de suavizado EMA (0 < alpha < 1).
        Valores bajos = más suave; valores altos = más reactivo.
    """

    def __init__(self, alpha: float = 0.1):
        self.alpha = alpha
        self._smoothed: float | None = None   # valor suavizado anterior

    def update(self, roi_gray: np.ndarray) -> int:
        """Devuelve la fila estimada de la línea del agua (coordenada y).

        Parámetros
        ----------
        roi_gray : np.ndarray
            Subimagen en escala de grises correspondiente a un cilindro.

        Retorna
        -------
        int
            Coordenada y de la línea del agua dentro de la ROI.
        """
        # Perfil de intensidad promedio por fila
        row_mean = roi_gray.mean(axis=1)          # shape: (H,)

        # Gradiente vertical: dónde cambia más la intensidad
        gradient = np.abs(np.diff(row_mean))      # shape: (H-1,)
        raw_y = float(np.argmax(gradient))

        # Primera llamada: inicializar con la estimación cruda
        if self._smoothed is None:
            self._smoothed = raw_y
        else:
            # Ecuación EMA: ŷ_t = alpha * y_t + (1 - alpha) * ŷ_{t-1}
            self._smoothed = (
                self.alpha * raw_y
                + (1.0 - self.alpha) * self._smoothed
            )

        return int(round(self._smoothed))
\end{lstlisting}

Los puntos clave del listado son:

\begin{itemize}
    \item \textbf{Líneas 19--20.} El perfil de intensidad por fila
    (\texttt{row\_mean}) reduce la imagen 2D a un vector 1D, eliminando
    variaciones horizontales irrelevantes.

    \item \textbf{Líneas 23--24.} \texttt{np.diff} calcula la diferencia
    entre filas consecutivas; \texttt{argmax} localiza la fila con el cambio
    más brusco, que corresponde a la transición agua--rata.

    \item \textbf{Líneas 27--33.} La actualización EMA corresponde
    directamente a la ecuación~\ref{eq:ema} del capítulo de marco teórico,
    con $\alpha = 0.1$ como valor por defecto.
\end{itemize}


% ════════════════════════════════════════════════════════════
%  FRAGMENTO 2 — YoloTracker
% ════════════════════════════════════════════════════════════

\subsubsection{Clase \texttt{YoloTracker}}

El Listado~\ref{lst:yolotracker} muestra la clase que integra YOLOv8 en modo
\textit{track} con ByteTrack. El diseño aísla la dependencia de Ultralytics
en un único punto: si la biblioteca no está disponible, el atributo
\texttt{available} indica al pipeline que debe usar el detector clásico.

\begin{lstlisting}[
  style=python,
  caption={Clase \texttt{YoloTracker} en \texttt{tracker.py}.},
  label={lst:yolotracker}
]
class YoloTracker:
    """Envuelve YOLOv8 en modo track (ByteTrack) para seguimiento de rata.

    Si Ultralytics no está instalado, ``available`` es False y el
    pipeline cae automáticamente al detector clásico.
    """

    WEIGHTS = Path("weights/rat.pt")
    TRACKER_CFG = "bytetrack.yaml"   # incluido en Ultralytics

    def __init__(self):
        self.available = False
        self._model = None
        try:
            from ultralytics import YOLO
            if not self.WEIGHTS.exists():
                raise FileNotFoundError(
                    f"Pesos no encontrados: {self.WEIGHTS}"
                )
            self._model = YOLO(str(self.WEIGHTS))
            self.available = True
        except Exception as exc:          # ImportError o FileNotFoundError
            warnings.warn(
                f"YoloTracker no disponible: {exc}. "
                "Se usará el detector clásico como respaldo."
            )

    def track(
        self,
        frame: np.ndarray,
        roi: tuple[int, int, int, int],
        conf: float = 0.25,
    ) -> list[dict]:
        """Ejecuta detección + seguimiento ByteTrack sobre una ROI.

        Parámetros
        ----------
        frame : np.ndarray
            Cuadro completo en BGR.
        roi : tuple[int, int, int, int]
            Región de interés (x, y, w, h) en píxeles del cuadro.
        conf : float
            Umbral mínimo de confianza para reportar una detección.

        Retorna
        -------
        list[dict]
            Lista de detecciones; cada elemento contiene:
            ``{'id': int, 'bbox': [x,y,w,h], 'conf': float}``.
            Lista vacía si no se detectó ningún objeto.
        """
        if not self.available:
            return []

        x, y, w, h = roi
        crop = frame[y:y+h, x:x+w]

        # persist=True mantiene el mismo ID entre cuadros consecutivos
        results = self._model.track(
            crop,
            tracker=self.TRACKER_CFG,
            persist=True,
            conf=conf,
            verbose=False,
        )

        detections = []
        if results and results[0].boxes is not None:
            for box in results[0].boxes:
                track_id = int(box.id) if box.id is not None else -1
                bx, by, bw, bh = box.xywh[0].tolist()
                detections.append({
                    "id":   track_id,
                    "bbox": [x + int(bx - bw/2),   # convertir a coords
                             y + int(by - bh/2),   # del cuadro completo
                             int(bw), int(bh)],
                    "conf": float(box.conf),
                })
        return detections
\end{lstlisting}

Los puntos clave del listado son:

\begin{itemize}
    \item \textbf{Líneas 11--23.} El bloque \texttt{try/except} aísla la
    dependencia de Ultralytics. Si la importación falla (por ejemplo, en un
    entorno sin GPU o sin la biblioteca instalada), el objeto queda en estado
    \texttt{available = False} sin lanzar una excepción que interrumpa el
    pipeline.

    \item \textbf{Líneas 41--43.} El cuadro completo se recorta a la ROI
    antes de pasarlo al detector. Esto reduce la carga computacional y evita
    que el modelo detecte ratas de otros cilindros dentro del mismo cuadro.

    \item \textbf{Línea 46.} \texttt{persist=True} es el parámetro que activa
    la continuidad de identidades entre cuadros en Ultralytics: ByteTrack
    recuerda qué ID asignó a cada trayectoria en el cuadro anterior y lo
    mantiene en el cuadro actual.

    \item \textbf{Líneas 55--62.} Las coordenadas del bounding box se
    devuelven en el sistema de referencia del cuadro completo (no de la ROI),
    para que el resto del pipeline pueda trabajar siempre en el mismo espacio
    de coordenadas.
\end{itemize}


% ════════════════════════════════════════════════════════════
%  FRAGMENTO 3 — Jerarquía de fallback en process_video()
% ════════════════════════════════════════════════════════════

\subsubsection{Jerarquía de respaldo en \texttt{process\_video()}}

El Listado~\ref{lst:fallback} muestra el núcleo del bucle de procesamiento
por cuadro dentro de \texttt{process\_video()}. Implementa la jerarquía de
tres niveles descrita en la sección de marco teórico: ByteTrack como
primera opción, detector clásico como respaldo y congelamiento de posición
como último recurso.

\begin{lstlisting}[
  style=python,
  caption={Bucle principal de seguimiento en \texttt{process\_video()}.
           Se omiten las partes de gating físico y escritura de resultados
           para centrarse en la jerarquía de respaldo.},
  label={lst:fallback}
]
for frame_idx, frame in enumerate(video_frames(cap)):

    for roi_id, roi in enumerate(rois):
        det = None

        # ── Nivel 1: ByteTrack (vía YoloTracker) ──────────────────
        if yolo_tracker.available:
            hits = yolo_tracker.track(frame, roi, conf=CONF_HIGH)
            if hits:
                det = hits[0]                 # exactamente 1 rata por ROI
                freeze_counter[roi_id] = 0   # reiniciar contador de freeze

        # ── Nivel 2: detector clásico (sustracción de fondo + Otsu) ─
        if det is None:
            bbox_classic = classic_detect(frame, roi, bg_model)
            if bbox_classic is not None and iou(
                bbox_classic, last_bbox[roi_id]
            ) > IOU_MIN:
                det = {"bbox": bbox_classic, "conf": CONF_CLASSIC,
                       "id": roi_id}
                freeze_counter[roi_id] = 0

        # ── Nivel 3: congelamiento con decaimiento de confianza ─────
        if det is None:
            freeze_counter[roi_id] += 1
            decay = CONF_FREEZE_INIT * (
                CONF_FREEZE_DECAY ** freeze_counter[roi_id]
            )
            det = {
                "bbox": last_bbox[roi_id],
                "conf": max(decay, CONF_FREEZE_MIN),
                "id":   roi_id,
            }

        # Actualizar estado y aplicar gating físico
        last_bbox[roi_id] = det["bbox"]
        waterline_y[roi_id] = wl_smoothers[roi_id].update(
            _crop_gray(frame, roi)
        )
        frame_results[roi_id].append(det)
\end{lstlisting}

Los puntos clave del listado son:

\begin{itemize}
    \item \textbf{Líneas 5--9.} Si \texttt{YoloTracker} está disponible y
    produce al menos una detección, se toma la primera (ya que hay exactamente
    una rata por ROI) y se reinicia el contador de \textit{freeze}.

    \item \textbf{Líneas 12--19.} El detector clásico solo se usa si
    ByteTrack no produjo resultado. La condición \texttt{iou > IOU\_MIN}
    actúa como \textit{gating} geográfico: si la detección clásica aparece
    demasiado lejos de la última posición conocida, se descarta como
    artefacto (reflejo o ruido morfológico).

    \item \textbf{Líneas 22--31.} El congelamiento asigna la última posición
    conocida con una confianza que decrece exponencialmente cada cuadro
    ($\texttt{conf} = c_0 \cdot \gamma^{k}$, donde $k$ es el número de
    cuadros congelados). Esto permite que el clasificador de conducta posterior
    descarte las predicciones de muy baja confianza, en lugar de tratarlas
    como detecciones válidas.

    \item \textbf{Líneas 34--38.} La línea del agua se actualiza en cada
    cuadro usando \texttt{WaterlineSmoother}, independientemente del nivel
    de respaldo que haya producido la detección.
\end{itemize}


% ══════════════════════════════════════════════════════════════
%  [G]  bib/referencias.bib  — SECCIÓN NUEVA
%
%  Las tres entradas de [A] deben quedar en el .bib.
%  Esta sección [G] es un recordatorio explícito de qué claves
%  se usan en cada archivo y con qué propósito, para facilitar
%  la revisión de consistencia antes de compilar.
%
%  Clave                     Usada en       Propósito
%  ─────────────────────────────────────────────────────────────
%  zhang2022_bytetrack       B, C, D, F     Paper principal ByteTrack ECCV 2022
%  bewley2016_sort           B, C           Paper SORT (comparación cuantitativa)
%  wojke2017_deepsort        B, C, D        Paper DeepSORT (comparación cuantitativa)
%  ultralytics_bytetrack_yaml C, D, F       Archivo bytetrack.yaml de Ultralytics
%
%  Verificar que las claves anteriores NO existan ya en
%  bib/referencias.bib con nombres distintos antes de insertar.
%  Las claves lukezic2017 y henriques2015 ya existen en el .bib
%  original — no duplicar.
%
%  Las entradas .bib están en la sección [A] de este archivo.
% ══════════════════════════════════════════════════════════════